\bookchapter{Глава 5: Пусть телефон понимает вас лучше — Базовые настройки}

\chapterintro{Потратьте несколько минут на настройку — и ощущения от использования изменятся кардинально. Эти настройки как «настройка инструмента» для телефона, чтобы он лучше соответствовал вашим привычкам.}

\booksection{Звук и вибрация}

Установите любимую мелодию звонка и уведомлений. Гибко используйте беззвучный режим (для совещаний) и режим «Не беспокоить» (для сна) — мир станет намного тише.

\begin{itemize}
  \item \textbf{Мелодия звонка}: выберите такую, которую вы услышите и которая не смущает.
  \item \textbf{Режим вибрации}: в беззвучном режиме оповещает вибрацией.
  \item \textbf{Режим «Не беспокоить»}: можно настроить исключения, например, пропускать звонки из адресной книги.
  \item \textbf{Разделение громкости}: громкость мультимедиа, звонка и будильника регулируются отдельно.
\end{itemize}

\booksection{Дисплей и яркость}

\begin{itemize}
  \item \textbf{Автояркость}: пусть телефон сам регулирует яркость в зависимости от освещения — удобно.
  \item \textbf{Защита зрения / Ночной режим}: фильтрует синий свет, глазам комфортнее вечером.
  \item \textbf{Тёмная тема}: экономит заряд (особенно эффективно для OLED-экранов) и не режет глаза ночью. Многие приложения тоже поддерживают тёмную тему.
  \item \textbf{Частота обновления экрана}: на некоторых Android-телефонах регулируется, чем выше (например, 120 Гц), тем плавнее прокрутка.
  \item \textbf{Автоблокировка}: установите 15 или 30 секунд — и безопасно, и экономит заряд.
  \item \textbf{Размер шрифта}: кажется, что буквы мелкие? Увеличьте.
\end{itemize}

\booksection{Обои, темы и рабочий стол}

Смените обои на красивые (статические или динамические) — и настроение уже другое. У брендовых Android-телефонов часто есть магазин тем, где можно полностью сменить иконки и стиль.

\textbf{Виджеты рабочего стола}: погода, календарь, управление музыкой на рабочем столе — доступ в одно касание. iOS и Android поддерживают, долгое нажатие на пустом месте рабочего стола для добавления.

\booksection{Настройки эффективности: превращаем телефон в ваш «внешний мозг»}

Эти функции из разряда «один раз настроил — и обратно не вернёшься», каждая из них заметно повышает вашу повседневную эффективность.

\booksubsection{Режим фокусировки / «Не беспокоить»: возвращаем внимание, украденное уведомлениями}

\begin{itemize}
  \item \textbf{Основная идея}: не «меньше смотреть в телефон», а «фокусироваться, когда нужно, и быть на связи, когда нужно».
  \item \textbf{Рекомендации по использованию}:\newline · Рабочий режим: пропускать только «избранные» контакты и необходимые приложения;\newline · Режим сна: автоматическое включение, только будильник;\newline · Фитнес / Чтение / Совещание: настройка под текущую деятельность.
  \item \textbf{Продвинутый уровень}: используйте автоматизацию + местоположение / время для автоматического переключения режимов. Например: iOS-команды (Shortcuts) позволяют настроить «автоматический переход в рабочий режим по прибытии в офис».
\end{itemize}

\booksubsection{Управление уведомлениями: чтобы красные точки не стали «бесконечным звонком»}

\begin{itemize}
  \item \textbf{Отключите маркетинговые уведомления}: 90\% push-уведомлений вам на самом деле не нужны.
  \item \textbf{Сортировка по важности}: закрепите коммуникационные (телефон, SMS, мессенджеры) наверху; новости, магазины, соцсети можно перевести в «тихие уведомления».
  \item \textbf{Группировка уведомлений}: сворачивайте однотипные уведомления вместе, чтобы одна разблокировка не топила вас в десятках сообщений.
\end{itemize}

\booksubsection{Быстрые жесты: опережайте на шаг}

\begin{itemize}
  \item \textbf{Скриншот}: кнопка питания + громкость вниз (iOS полноэкранный) / питание + громкость вниз (Android универсально).
  \item \textbf{Длинный скриншот / Прокручиваемый скриншот}: после обычного скриншота нажмите опцию «Длинный скриншот».
  \item \textbf{Быстрое сканирование QR}: добавьте иконку «Сканер» в центр управления / быстрые настройки.
  \item \textbf{Касание задней панели (iPhone)}: в [Настройки] > [Универсальный доступ] > [Касание] > [Касание задней панели] назначьте действия вроде «Скриншот» или «Блокировка экрана».
  \item \textbf{Голосовая активация}: Siri / XiaoAi / XiaoYi / XiaoBu — окликните, чтобы узнать погоду или поставить будильник.
\end{itemize}

\booksubsection{Клавиатура: самый недооценённый инструмент эффективности}

\begin{itemize}
  \item \textbf{Выберите правильную клавиатуру}:\newline · Sogou / Baidu: большой словарь, быстрое обновление популярных слов;\newline · Xunfei: лучшая точность голосового ввода, идеально для конспектов совещаний;\newline · Gboard (Google): лаконичная, сильная кроссплатформенная синхронизация, поддержка свайп-ввода;\newline · Встроенная iOS: лучшая конфиденциальность, глубокая интеграция с системой.
  \item \textbf{Обязательные переключатели}: интеллектуальное исправление, подсказки слов, синхронизация облачного словаря, ввод цифр и символов смахиванием вверх.
  \item \textbf{Совет по конфиденциальности}: осторожнее с функциями «подсказки целых предложений» или «персонализированные рекомендации» — они могут отправлять вводимый текст.
\end{itemize}

\booksubsection{Рабочий стол и многозадачность: пусть рабочий процесс течёт свободно}

\begin{itemize}
  \item \textbf{Разделённый экран / Картинка в картинке}: смотрите видео и отвечайте на сообщения, выходите в интернет и делайте заметки.
  \item \textbf{Маленькое окно / Плавающее окно}: WeChat / Douyin поддерживают «режим маленького окна» — ответил и свернул.
  \item \textbf{Закрепление частых приложений}: в Dock-панели (внизу экрана) держите только 4 самых частых приложения, чтобы сократить время поиска.
\end{itemize}

\begin{tcbinfo}{Чек-лист «Первого рабочего дня»}
При получении нового телефона или после крупного обновления системы рекомендуем потратить 10 минут на эти настройки:\newline (1) Пароль блокировки + биометрия (2) Группировка уведомлений + отключение маркетинговых push (3) Размер шрифта / частота обновления / тёмная тема\newline (4) Клавиатура по умолчанию / браузер / почта (5) Экстренные контакты / медицинская карта (6) Настройки резервного копирования

\end{tcbinfo}

\booksection{Специальные возможности и универсальный доступ: сила технологий во благо}

Не пугайтесь названия «Специальные возможности» — эти функции доступны каждому. Изначально они создавались для людей с ограниченными возможностями, но полезны и обычным пользователям — временно injured рука с ребёнком на руках, сломались наушники в метро и ничего не слышно, надели очки для чтения и не видите мелкий текст… всё это может выручить. По сути, универсальный дизайн воплощает идею\textbf{технологий во благо}: дать каждому возможность равноправно наслаждаться цифровой жизнью.

\booksubsection{«Глаза» для слабовидящих — полное руководство по экранным дикторам}

Экранный диктор — основной инструмент для незрячих пользователей. Он\textbf{озвучивает всё видимое содержимое экрана}, а пользователь управляет телефоном жестами и голосом. Ниже описаны методы использования iOS VoiceOver и Android TalkBack.

\begin{tcbinfo}{iOS VoiceOver (экранный диктор)}
\textbf{Путь включения}: [Настройки] > [Универсальный доступ] > [VoiceOver], включите переключатель. Также можно быстро включить/выключить тройным нажатием боковой кнопки (или кнопки «Домой»).

\textbf{Основные жесты}:

\begin{enumerate}
  \item Смахивание одним пальцем: поэлементная навигация по экрану;
  \item Смахивание двумя пальцами: прокрутка страницы вверх/вниз и влево/вправо;
  \item Двойное касание одним пальцем: активация текущего элемента (аналог «нажатия»);
  \item Смахивание тремя пальцами: переключение страниц;
  \item Тройное касание одним пальцем: увеличение или выбор (отличается от обычного тройного касания).
\end{enumerate}

\textbf{Персонализация}: [VoiceOver] > [Речь], можно настроить\textbf{скорость озвучивания}(по умолчанию около 50\%, можно увеличить до 70–80\% для более плавного восприятия),\textbf{высоту тона}, переключить голосовые пакеты. iOS предлагает «Улучшенную речь» и «Голосовой контроллер» — первая уменьшает дыхательные шумы и паузы, делая звучание более естественным.

\end{tcbinfo}

\begin{tcbinfo}{Android TalkBack}
\textbf{Путь включения}: [Настройки] > [Специальные возможности] > [TalkBack], включите переключатель. Или зажмите кнопки громкости вверх и вниз на 3 секунды для быстрого включения (поддерживается не всеми брендами).

\textbf{Основные жесты}:

\begin{enumerate}
  \item Смахивание одним пальцем влево-вправо: переключение фокуса между элементами;
  \item Двойное касание одним пальцем: активация элемента;
  \item Смахивание двумя пальцами вверх/вниз и влево/вправо: прокрутка;
  \item Одинарное касание тремя пальцами: вход в режим чтения (непрерывное озвучивание);
  \item Одинарное касание четырьмя пальцами: пауза/продолжение озвучивания.
\end{enumerate}

\textbf{Персонализация}: [TalkBack] > [Настройки синтеза речи], можно настроить\textbf{скорость},\textbf{высоту тона}, переключить\textbf{движок синтеза речи}(рекомендуется установить Google Text-to-Speech для более естественного озвучивания).

\end{tcbinfo}

\booksubsection{Делаем мелкое крупным, размытое чётким — настройка визуального отображения}

\booksubsection{Жест увеличения}

\begin{itemize}
  \item \textbf{iOS}: [Настройки] > [Универсальный доступ] > [Увеличение], после включения\textbf{тройное касание экрана}увеличивает его. По умолчанию — полноэкранное увеличение, можно выбрать оконное (локальное увеличительное окно).\textbf{Разведение двух пальцев регулирует кратность (от 1,5x до 15x)}, двойное касание тремя пальцами фиксирует текущий масштаб.
  \item \textbf{Android}: [Настройки] > [Специальные возможности] > [Увеличение], тройное касание увеличивает, повторное тройное касание выходит; или\textbf{тройное касание экрана трижды}для временного увеличения (до 8 крат).
\end{itemize}

\booksubsection{Высококонтрастный текст / Жирный текст}

\begin{itemize}
  \item \textbf{iOS}: [Настройки] > [Универсальный доступ] > [Дисплей и размер текста], включите «Жирный текст», «Повышенная контрастность», «Понижение прозрачности» и другие переключатели. «Умная инверсия» сохраняет текст чёрным на светлом фоне, избегая белого текста на белом фоне.
  \item \textbf{Android}: [Настройки] > [Дисплей] > [Размер шрифта], можно выбрать «Мелкий, Обычный, Крупный, Очень крупный». Huawei/Xiaomi дополнительно предлагают переключатели «Жирный» и «Высококонтрастный текст».
\end{itemize}

\booksubsection{Цветовые фильтры (для дальтоников)}

\begin{itemize}
  \item \textbf{iOS}: [Настройки] > [Универсальный доступ] > [Дисплей и размер текста] > [Цветовые фильтры], предлагает\textbf{четыре режима для дальтоников: красно-зелёный/зелёно-красный/сине-жёлтый}, а также «Оттенки серого» и «Инверсия цвета». Рекомендуется сначала пройти тест на цветовосприятие (встроенный в систему), затем выбрать соответствующий фильтр.
  \item \textbf{Android}: [Настройки] > [Специальные возможности] > [Коррекция цвета], включает\textbf{три режима для дальтоников}+ переключатель «Инверсия цвета».
\end{itemize}

\booksubsection{AI-описание изображений: пусть AI станет вашими «глазами»}

Когда экранный диктор встречает изображение, традиционно он либо пропускает его, либо произносит лишь «изображение». Но теперь AI может\textbf{описать содержание изображения одним предложением}, что очень помогает при распознавании капчи, чтении картинок и меню.

\begin{tcbinfo}{Поддержка AI-описания изображений в разных системах}
\begin{itemize}
  \item \textbf{iOS 15+}: в [Настройки] > [Универсальный доступ] > [VoiceOver] > [Описание речи] включите «Описание изображений». Система использует\textbf{локальную AI-модель}для автоматической генерации текстового описания изображений (требуется iPhone 12 и новее).
  \item \textbf{Новая версия TalkBack}: Google TalkBack в новых версиях интегрирует функцию\textbf{вопросов по изображению}, можно спросить «Что на этой картинке?», и AI ответит в реальном времени (не зависит строго от версии Android, достаточно обновить приложение TalkBack).
  \item \textbf{Сторонние приложения}: «Seeing AI» (iOS), «Lookout» (Android), «XiaoYi Bang Bang» (Huawei) — отличные AI-визуальные помощники, распознают предметы, читают документы, определяют выражения лиц.
\end{itemize}

\end{tcbinfo}

\begin{tcbpractice}{Сценарии использования}
Попробуйте спросить у уличной фотографии: «Сколько людей на этой картинке? Что они делают?» или откройте приложение доставки еды, чтобы AI прочитал названия блюд из меню —\textbf{это по-настоящему «технологии, делающие жизнь равнее»}.

\end{tcbpractice}

\booksubsection{«Уши» для слабослышащих — полное руководство по живым субтитрам}

Живые субтитры (Live Captions) — это\textbf{системная функция преобразования речи в текст}, особенно полезная для слабослышащих пользователей в шумной обстановке, библиотеке или при иностранной речи.

\begin{tcbinfo}{Включение и оптимизация живых субтитров}
\textbf{iOS}: [Настройки] > [Универсальный доступ] > [Живые субтитры], включите переключатель. Система будет определять весь звук с устройства (звонки, видео, подкасты, FaceTime, голосовые заметки) и автоматически отображать субтитры в верхней части экрана.

\textbf{Android 12+}: нажмите\textbf{кнопку громкости}нажмите значок «Живые субтитры» во всплывающем окне; или включите в [Настройки] > [Специальные возможности] > [Живые субтитры].

\textbf{Советы по повышению точности}:

\begin{enumerate}
  \item По возможности используйте в\textbf{тихой обстановке}, вдали от кондиционеров и гудящих машин;
  \item Наивысшая точность достигается при\textbf{чётком произношении и умеренном темпе речи};
  \item Разговор нескольких человек снижает точность, можно использовать режим «Показывать только одного»;
  \item Поддерживаются основные языки: путунхуа, кантонский, английский,\textbf{точность для редких языков ниже};
  \item При фоновой музыке или посторонних голосах распознавание будет прерывистым.
\end{enumerate}

\end{tcbinfo}

\booksubsection{Двухканальный режим звонков: речь в текст + текст в речь}

Когда пользователь\textbf{и не слышит, и не может говорить}(например, внезапная потеря голоса, восстановление после операции на гортани, сильное заикание), телефон может стать «переводчиком».

\begin{tcbinfo}{Трансляция разговора в реальном времени (TTY / RTT)}
\textbf{iOS}: [Настройки] > [Универсальный доступ] > [RTT/TTY], после включения во время звонка на экране появляется клавиатура, речь собеседника в реальном времени преобразуется в текст,\textbf{а набранный вами текст озвучивается}для собеседника.

\textbf{Android}: [Настройки] > [Специальные возможности] > [Живая транскрипция] + [Звуковые уведомления]. «Живая транскрипция» отображает речь собеседника в реальном времени; «Звуковые уведомления» при невозможности говорить преобразуют ваш ввод в голос.

\textbf{Совместимость}:

\begin{itemize}
  \item \textbf{RTT (Real-Time Text)}: требуется поддержка обеими сторонами, текст отображается по мере ввода, без задержки, основные операторы уже поддерживают;
  \item \textbf{TTY (телетайп)}: требуется специальное TTY-устройство, постепенно заменяется RTT;
  \item Функция\textbf{«Голос в текст»}в WeChat/QQ также является альтернативой: долгое нажатие на голосовое сообщение «В текст».
\end{itemize}

\end{tcbinfo}

\booksubsection{«Маленькие, но практичные» функции для всех}

\begin{itemize}
  \item \textbf{Касание задней панели}(iOS 14+): [Настройки] > [Универсальный доступ] > [Касание] > [Касание задней панели], можно назначить\textbf{двойное/тройное касание}для скриншота, блокировки, центра управления и других действий. Особенно удобно, когда руки заняты посылкой.
  \item \textbf{AssistiveTouch / Плавающая кнопка}: iOS [Универсальный доступ] > [Касание] > [AssistiveTouch], Android [Специальные возможности] > [Плавающая навигация]. Виртуальная кнопка заменяет жесты или физические клавиши — выручает при неработающих кнопках или управлении одной рукой.
  \item \textbf{Управление переключением}: управление телефоном через внешние переключатели, одиночные нажатия, движения головы —\textbf{разработано для пользователей с нарушениями опорно-двигательного аппарата}. iOS [Универсальный доступ] > [Управление переключением], Android [Специальные возможности] > [Управление переключателями].
  \item \textbf{Распознавание звуков}(iOS): [Универсальный доступ] > [Распознавание звуков], после включения телефон «слушает» дверной звонок, плач ребёнка, пожарную сигнализацию, сирену и другие звуки,\textbf{при обнаружении — вибрация + уведомление на экране}, очень ценно для слабослышащих пользователей и во время сна.
  \item \textbf{LED-вспышка для оповещений}: iOS [Универсальный доступ] > [Аудио/Видео] > [LED-вспышка для оповещений], Android [Настройки] > [Уведомления] > [Оповещение вспышкой]. При звонках и уведомлениях вспышка мигает —\textbf{в шумной обстановке можно воспринимать без звука}.
  \item \textbf{Монофонический звук}: [Настройки] > [Универсальный доступ] > [Аудио] > [Монофонический звук], объединяет левый и правый каналы. Пользователи с односторонним нарушением слуха не теряют половину звуковой информации.
\end{itemize}

\begin{tcbinfo}{Философия универсального дизайна Apple}
Apple вкладывает огромные ресурсы в сферу универсального доступа, ежегодно обновляя и добавляя новые функции. На официальном сайте Apple есть специальная страница универсального доступа с подробным описанием всех функций. Функции универсального доступа iOS обычно более проработаны и плавны, чем на Android, поэтому многие пользователи с ограниченными возможностями выбирают Apple.

\end{tcbinfo}

\begin{tcbpractice}{Потратьте 30 секунд, чтобы попробовать}
Найдите свободное время и попробуйте следующее:

\begin{enumerate}
  \item Включите\textbf{VoiceOver / TalkBack}и попробуйте, не глядя на экран, пролистать WeChat тремя пальцами;
  \item Включите\textbf{живые субтитры}и поговорите в телефон, проверяя точность распознавания;
  \item Увеличьте фрагмент текста до 200\%\textbf{тройным касанием}и испытайте повседневность слабовидящего пользователя;
  \item Отключите все эти функции и подумайте: «Если бы мне всегда были нужны эти функции, как бы я хотел, чтобы они были спроектированы?»
\end{enumerate}

\end{tcbpractice}
